\documentclass[11pt,a4paper]{article}

% --- PACCHETTI BASE E LINGUA ---
\usepackage[utf8]{inputenc}
\usepackage[T1]{fontenc}
\usepackage[italian]{babel}
\usepackage{lmodern}
\usepackage{microtype}

% --- LAYOUT E GEOMETRIA ---
\usepackage{geometry}
\geometry{a4paper, top=2.5cm, bottom=2.5cm, left=2.2cm, right=2.2cm}
\usepackage{setspace}
\onehalfspacing

% --- GRAFICA E TABELLE ---
\usepackage{booktabs}
\usepackage{tabularx}
\usepackage{amsmath}
\usepackage{siunitx}
\sisetup{output-decimal-marker={,}, per-mode=symbol}
\usepackage{pgfplots}
\pgfplotsset{compat=1.18}
\usepackage{float}
\usepackage{caption}
\captionsetup{font=small, labelfont=bf}
\usepackage[dvipsnames]{xcolor}
\usepackage{footmisc}
\usepackage[hidelinks]{hyperref}

% --- STILE TITOLI ---
\usepackage{titlesec}
\titleformat{\section}{\normalfont\Large\bfseries}{\thesection}{0.8em}{}
\titleformat{\subsection}{\normalfont\large\bfseries}{\thesubsection}{0.8em}{}

% --- COLONNE TABULARX ---
\newcolumntype{Y}{>{\raggedright\arraybackslash}X}
\newcolumntype{C}{>{\centering\arraybackslash}X}
\newcommand{\gradi}{$^{\circ}$C}

\title{\vspace{-1.4cm}\textbf{IL CLIMA CAMBIA ANCHE A CASA NOSTRA}\\[4pt]
\large Quanto fa più caldo nei 32 comuni P.A.E.S.C. e come questo impatta la nostra vita quotidiana}
\author{A cura del GAL Versante Laziale del PNALM}
\date{Aggiornamento Dati: Giugno 2026}

\begin{document}

\maketitle

\section*{Sommario}
Spesso pensiamo al riscaldamento globale come a un problema lontano, che riguarda i ghiacciai dell'Artico o gli orsi polari. Eppure le estati sono sempre più roventi e le piogge sempre più imprevedibili anche fuori dalla nostra porta di casa. Questo documento, a supporto di una nuova piattaforma web di sensibilizzazione, traduce i complessi dati climatici satellitari europei nella realtà quotidiana dei 32 comuni del GAL Verla. Il nostro obiettivo non è fare allarmismo, ma rendere visibile e comprensibile a tutti un fenomeno che sta lentamente mutando il nostro territorio, le nostre abitudini e la nostra economia. I numeri che seguono raccontano una storia semplice e netta: rispetto al clima dei nostri nonni, qui fa stabilmente più caldo di oltre due gradi, e d'estate anche di più.

\section{Perché guardare al clima del proprio paese}
Quando si parla di cambiamento climatico, la domanda che sorge spontanea in ognuno di noi è: ``Sì, ma da me quanto fa più caldo rispetto a una volta?''. Non è una semplice curiosità, ma una questione vitale per un territorio come quello del GAL Versante Laziale del Parco Nazionale d'Abruzzo, Lazio e Molise (GAL Verla). Parliamo di 32 amministrazioni comunali in provincia di Frosinone, estese su oltre \SI{1168}{\kilo\meter\squared}, in gran parte classificate come zona montana. In quest'area l'agricoltura, il turismo e la qualità della vita dipendono intimamente dall'equilibrio del clima locale: una vendemmia anticipata, un'estate più lunga, una sorgente più povera d'acqua non sono dettagli, ma fatti che incidono sull'economia e sulla salute delle persone.

Le ondate di calore fanno notizia perché sono visibili e improvvise, ma sono solo la punta dell'iceberg. Il segnale che conta davvero è lo spostamento lento e silenzioso della \textbf{media termica}: è la media a dirci se un territorio sta cambiando pelle. Dagli anni Ottanta il bacino del Mediterraneo si scalda a un ritmo superiore del 20\% rispetto alla media globale, e sono proprio le zone interne e d'altura --- come gran parte del nostro comprensorio --- a riscaldarsi più in fretta. Per capire l'entità del fenomeno a casa nostra abbiamo unito il rigore della scienza alla chiarezza della divulgazione.

\section{Da dove vengono i dati: i satelliti europei}
Per calcolare l'aumento delle temperature non ci siamo affidati alle singole centraline meteo locali, che spesso non esistono nei piccoli comuni montani o forniscono dati falsati dall'asfalto cittadino. Abbiamo invece utilizzato la \textbf{rianalisi climatica}, ovvero il potente database \textit{ERA5} (e la sua versione ad alta risoluzione \textit{ERA5-Land}) del programma europeo \textbf{Copernicus}.

Immaginate un supercomputer che incrocia i dati dei satelliti con migliaia di osservazioni al suolo, creando una mappa termica oggettiva del nostro territorio, ora dopo ora, dal 1961 fino a oggi. Poiché in montagna fa ovviamente più freddo che a valle, i dati sono stati elaborati tenendo conto dell'altitudine esatta di ogni singolo comune, applicando un correttivo matematico noto come \textbf{downscaling altimetrico}. Sapendo che la temperatura scende di circa \SI{6.5}{\celsius} ogni mille metri di dislivello, l'algoritmo calcola la temperatura giusta per ogni paese:
\begin{equation*}
T_{\text{comune}} = T_{\text{satellite}} - 6{,}5 \cdot \Delta z \quad (\Delta z \text{ in km}).
\end{equation*}

Il risultato è un confronto netto e inequivocabile fra le temperature del passato e quelle di oggi. Come ``passato'' usiamo il periodo di riferimento ufficiale stabilito dall'Organizzazione Meteorologica Mondiale, la \textbf{media degli anni 1961--1990}: è la fotografia del clima che hanno conosciuto i nostri genitori e nonni. Avendo recuperato l'intera serie storica, oggi possiamo confrontarla con gli ultimi anni (2011--2025) senza buchi e senza forzature.

\section{L'App Web: uno strumento per capire e sensibilizzare}
Per rendere questi dati accessibili a chiunque abbiamo creato l'applicativo web \url{https://clima-galverla.vercel.app}. Questa piattaforma non è pensata per i soli scienziati, ma nasce come un vero strumento di \textbf{sensibilizzazione civica}: si sceglie il proprio comune e, in un istante, si legge quanto è cambiato il clima dove si vive.

Dire che la temperatura media è aumentata di ``due gradi'' può sembrare poco. Per questo l'app funge da \textbf{traduttore}: prende i gradi centigradi e li trasforma in impatti concreti sulla nostra vita di tutti i giorni, riassunti nella Tabella~\ref{tab:impatti}. La scienza ci dice infatti che l'innalzamento delle temperature porta conseguenze misurabili e quotidiane.

\begin{table}[H]
\centering
\caption{Come il clima impatta la nostra vita quotidiana (i fattori di traduzione dell'app).}
\label{tab:impatti}
\renewcommand{\arraystretch}{1.35}
\begin{tabularx}{\textwidth}{Y C Y}
\toprule
\textbf{Ambito della vita} & \textbf{L'impatto misurato} & \textbf{Significato pratico} \\
\midrule
Salute e riposo & $-12$ minuti di sonno & Per ogni ``notte tropicale'' dormiamo meno e peggio, e ci svegliamo più stanchi e vulnerabili. \\
Bollette e consumi & $+13\%$ di energia elettrica per grado & Rinfrescare le case con i condizionatori costa di più, ogni estate che passa. \\
Agricoltura e acqua & $+6\%$ di fabbisogno irriguo per grado & Le campagne hanno bisogno di molta più acqua, proprio mentre le siccità si allungano. \\
Cibo e raccolti & fino a $-5\%$ di resa dei cereali per grado & Meno produzione e vendemmie anticipate: il caldo cambia anche quello che mettiamo in tavola. \\
\bottomrule
\end{tabularx}
\end{table}

Oltre al ``traduttore'', l'app offre alcuni strumenti pensati per coinvolgere chi la usa: la \textbf{striscia degli anni}, che con un colpo d'occhio mostra il passaggio dal blu (freddo) al rosso (caldo) anno per anno; la funzione \textbf{``che tempo faceva il giorno della tua nascita''}, che confronta la propria data di nascita (dal 1961) con lo stesso giorno di oggi; la \textbf{classifica} dei comuni più colpiti dal caldo notturno; uno \textbf{scenario al 2050}; e una \textbf{card da condividere} sui social, per far girare il dato del proprio paese.

\section{I risultati sul nostro territorio}
Guardando ai numeri dei 32 comuni, il quadro è chiarissimo: \textbf{le nostre estati si stanno infiammando} e gli inverni si addolciscono. Rispetto al clima del 1961--1990, la temperatura media dell'intero distretto è salita di circa \textbf{$+2{,}3$~\gradi}, con un picco estivo che supera i tre gradi.

\subsection{L'estate rovente: il caso di Atina}
Prendiamo come esempio il comune di Atina (\SI{490}{\meter} s.l.m.). La temperatura media annua è cresciuta di circa \SI{2.3}{\celsius}, ma è nei mesi caldi che si registra lo scarto maggiore (Figura~\ref{fig:atina}): a giugno e ad agosto si tocca un impressionante \textbf{$+3{,}2$~\gradi} rispetto al passato. Questo significa ondate di calore più intense e durature, capaci di mettere a dura prova le fasce più deboli della popolazione, e notti in cui diventa difficile riposare.

\begin{figure}[H]
\centering
\begin{tikzpicture}
\begin{axis}[
    width=0.92\textwidth, height=7cm,
    ybar, bar width=16pt,
    title={\textbf{Scarto termico mese per mese ad Atina: quanto fa più caldo rispetto agli anni '60-'80}},
    ylabel={Gradi di differenza ($^{\circ}$C)},
    ymin=0, ymax=4,
    xtick=data,
    xticklabels={Gen,Feb,Mar,Apr,Mag,Giu,Lug,Ago,Set,Ott,Nov,Dic},
    nodes near coords, nodes near coords align={vertical},
    every node near coord/.append style={font=\scriptsize},
    grid=major, major grid style={dashed, gray!30},
    axis x line=bottom, axis y line=left,
    enlarge x limits=0.06,
]
\addplot[fill=red!65, draw=black!70] coordinates {
    (1,1.9)(2,2.1)(3,2.1)(4,2.4)(5,1.6)(6,3.2)(7,3.0)(8,3.2)(9,2.0)(10,1.9)(11,2.3)(12,1.9)
};
\addplot[sharp plot, no markers, black, dashed, thick] coordinates {(0.4,2.3)(12.6,2.3)};
\end{axis}
\end{tikzpicture}
\caption{L'istogramma mostra di quanti gradi sono aumentate le temperature ad Atina, mese per mese, rispetto alla media del periodo 1961--1990. La linea tratteggiata è l'aumento medio annuo ($+2{,}3$~\gradi). Tutti i mesi sono più caldi del passato; il picco estivo è evidente.}
\label{fig:atina}
\end{figure}

\subsection{Il disagio notturno e il rifugio della montagna}
Il vero indicatore del disagio fisico sono le cosiddette \textbf{``notti tropicali''}, quelle notti afose in cui il termometro non scende mai sotto i \SI{20}{\celsius}. Come illustrato nella Figura~\ref{fig:notti}, qui entra in gioco la geografia del nostro territorio. Nei comuni di fondovalle le notti tropicali sono letteralmente esplose, passando da una decina a oltre quaranta notti insonni all'anno. Man mano che si sale di quota, però, la montagna offre ancora un rifugio: oltre i 700 metri l'aria frizzante notturna resiste e le notti tropicali restano vicine a zero. È la stessa ragione per cui due paesi vicini, ma a quote diverse, vivono estati molto differenti.

\begin{figure}[H]
\centering
\begin{tikzpicture}
\begin{axis}[
    width=0.92\textwidth, height=7cm,
    title={\textbf{Notti insonni (notti tropicali) e altitudine: chi soffre di più il caldo?}},
    xlabel={Altitudine del comune (metri sul livello del mare)},
    ylabel={Notti oltre i $20\,^{\circ}$C all'anno},
    xmin=0, xmax=1100, ymin=-2, ymax=55,
    grid=major, major grid style={dashed, gray!30},
    only marks, mark size=2.4pt,
]
\addplot+[only marks, mark=*, mark options={fill=Maroon, draw=black!60}] coordinates {
    (60,47)(110,43)(130,42)(145,26)(188,20)(235,30)(250,29)(270,38)(350,25)
    (412,9)(460,0)(490,4)(550,0)(650,0)(725,0)(850,0)(926,0)(1050,0)
};
\end{axis}
\end{tikzpicture}
\caption{Ogni punto è un comune del GAL Verla. Più un paese si trova in basso (a sinistra), maggiore è il numero di notti afose (in alto). Chi vive in collina o in montagna gode ancora di notti fresche.}
\label{fig:notti}
\end{figure}

\subsection{La mappa termica dei nostri comuni}
La Tabella~\ref{tab:comuni} offre una fotografia chiara e diretta per alcuni dei comuni più rappresentativi. I dati mostrano come, a fronte di un innalzamento medio costante (attorno ai due gradi e mezzo ovunque), le conseguenze sul caldo notturno varino fortemente con la posizione geografica.

\begin{table}[H]
\centering
\caption{Il clima che cambia: aumento termico e notti afose nei comuni rappresentativi.}
\label{tab:comuni}
\renewcommand{\arraystretch}{1.25}
\begin{tabularx}{\textwidth}{Y C C C C}
\toprule
\textbf{Comune} & \textbf{Altitudine} & \textbf{Aumento annuo} & \textbf{Picco estivo} & \textbf{Notti tropicali (passato $\rightarrow$ oggi)} \\
\midrule
Pontecorvo & \SI{60}{\meter} & $+2{,}3$~\gradi & $+3{,}2$~\gradi & $20 \rightarrow 47$ \\
Piedimonte San Germano & \SI{130}{\meter} & $+2{,}3$~\gradi & $+3{,}2$~\gradi & $16 \rightarrow 42$ \\
Roccasecca & \SI{235}{\meter} & $+2{,}3$~\gradi & $+3{,}2$~\gradi & $11 \rightarrow 30$ \\
Ripi & \SI{350}{\meter} & $+2{,}3$~\gradi & $+3{,}2$~\gradi & $9 \rightarrow 25$ \\
Atina & \SI{490}{\meter} & $+2{,}3$~\gradi & $+3{,}2$~\gradi & $1 \rightarrow 4$ \\
Campoli Appennino & \SI{650}{\meter} & $+2{,}4$~\gradi & $+3{,}2$~\gradi & $0 \rightarrow 0$ \\
San Donato Val di Comino & \SI{728}{\meter} & $+2{,}4$~\gradi & $+3{,}2$~\gradi & $0 \rightarrow 0$ \\
Terelle & \SI{1050}{\meter} & $+2{,}3$~\gradi & $+3{,}2$~\gradi & $0 \rightarrow 0$ \\
\bottomrule
\end{tabularx}
\end{table}

\subsection{Uno sguardo al 2050: se non cambia nulla}
I numeri visti finora raccontano ciò che è già accaduto. Ma la stessa app prova anche a guardare avanti. Proiettando il ritmo di riscaldamento osservato, entro il \textbf{2050} il distretto potrebbe scaldarsi di un altro grado e mezzo abbondante rispetto a oggi. Per i fondovalle significherebbe un salto drammatico nelle notti afose: a Pontecorvo, per esempio, si passerebbe dalle attuali \num{47} a quasi \num{90} notti tropicali all'anno --- praticamente l'intera estate trascorsa con oltre venti gradi anche di notte. Non è una profezia, ma uno scenario di avvertimento: dipende da quanto sapremo ridurre, da qui in avanti, le emissioni che alimentano il riscaldamento.

\section{La consapevolezza è il primo passo}
I dati elaborati in questa relazione, e resi accessibili tramite la piattaforma web, dimostrano che il clima \textbf{è già cambiato}. Portare questi numeri sui nostri smartphone non serve a spaventarci, ma a prepararci.

Capire che dovremo ripensare i nostri consumi idrici, proteggere le persone anziane durante i mesi estivi, isolare meglio le abitazioni e accompagnare l'agricoltura verso colture e pratiche più resistenti al caldo è il primo passo per costruire un territorio capace di adattarsi. Sono le stesse direzioni indicate dal Piano d'Azione per l'Energia Sostenibile e il Clima (PAESC) del distretto: comunità energetiche rinnovabili, efficientamento degli edifici, forestazione, gestione attenta dell'acqua. Il riscaldamento globale non è più invisibile: grazie ai dati scientifici ora possiamo leggerlo, capirlo e, soprattutto, affrontarlo insieme.

\vspace{0.4cm}
\noindent\footnotesize\textit{Nota sui dati.} Elaborazione su dati ERA5 / ERA5-Land del Copernicus Climate Change Service (ECMWF). Serie osservata 1961--2025, confronto rispetto alla normale climatica WMO 1961--1990. Le medie del passato sono ricavate dalle temperature minime e massime giornaliere; gli impatti quotidiani usano fattori medi tratti da fonti ufficiali (ISPRA, EIA, Copernicus/HESS, CMCC, IPCC, MedECC). Strumento a fini divulgativi.

\end{document}
